\documentclass[11pt,a4paper]{article}
\usepackage[margin=1in]{geometry}
\usepackage{booktabs}
\usepackage{tabularx}
\usepackage{longtable}
\usepackage{enumitem}
\usepackage{fancyhdr}
\usepackage{hyperref}
\usepackage{xcolor}
\usepackage{tcolorbox}
\usepackage{amssymb}

\definecolor{headerblue}{HTML}{1E40AF}
\definecolor{shieldgreen}{HTML}{166534}
\definecolor{layergray}{HTML}{F3F4F6}

\hypersetup{colorlinks=true,linkcolor=headerblue,urlcolor=headerblue}

\pagestyle{fancy}
\fancyhf{}
\fancyhead[L]{\small\textbf{VajraKavach — Attack Mitigation Report}}
\fancyhead[R]{\small\textbf{CONFIDENTIAL}}
\fancyfoot[C]{\thepage}
\renewcommand{\headrulewidth}{0.5pt}

\title{
  \vspace{-1cm}
  \textbf{\Huge VajraKavach}\\[0.3cm]
  \textbf{\Large Attack Mitigation Report}\\[0.3cm]
  \large Preventing Attacks from Succeeding\\[0.8cm]
  \normalsize
  \begin{tabular}{ll}
    \textbf{Report Date:} & April 7, 2026 \\
    \textbf{Classification:} & CONFIDENTIAL \\
    \textbf{Prepared By:} & VajraKavach Security Team \\
    \textbf{Platform:} & Secure Job Search \& Professional Networking \\
    \textbf{Target:} & \texttt{https://192.168.3.44} \\
    \textbf{Course:} & CSE 345/545, IIIT-Delhi \\
  \end{tabular}
}
\author{}
\date{}

\begin{document}
\maketitle
\thispagestyle{fancy}

\newpage
\tableofcontents
\newpage

% ================================================================
\section{Executive Summary}
% ================================================================

Attack mitigation refers to the preventive security controls that stop attacks from succeeding in the first place. VajraKavach implements \textbf{23 distinct mitigation mechanisms} across four architectural layers: kernel/OS, reverse proxy (Nginx), application (FastAPI), and database (PostgreSQL).

This report catalogs every mitigation mechanism, explains its purpose, identifies the implementation location, and assesses its effectiveness. The platform follows a \textbf{defense-in-depth} strategy where multiple overlapping controls ensure that the failure of any single mechanism does not result in a successful attack.

% ================================================================
\section{Network \& Transport Layer Mitigation}
% ================================================================

\subsection{HTTPS Enforcement}

\begin{tcolorbox}[colback=shieldgreen!5, colframe=shieldgreen]
\textbf{Mechanism:} All HTTP (port 80) traffic is rejected with an immediate connection close (HTTP 444).\\
\textbf{Location:} \texttt{nginx/nginx.conf} — \texttt{return 444}\\
\textbf{Purpose:} Prevents credential leakage during redirect. Unlike a 301 redirect, 444 closes the TCP connection without sending any response body, eliminating the risk of transmitting sensitive data over plaintext.\\
\textbf{Effectiveness:} \textcolor{shieldgreen}{\textbf{STRONG}} — No HTTP traffic is served under any circumstance.
\end{tcolorbox}

\subsection{TLS Configuration}

\begin{tcolorbox}[colback=shieldgreen!5, colframe=shieldgreen]
\textbf{Mechanism:} TLS 1.2 and TLS 1.3 only, with AEAD cipher suites.\\
\textbf{Location:} \texttt{nginx/nginx.conf}\\
\textbf{Configuration:}
\begin{verbatim}
ssl_protocols TLSv1.2 TLSv1.3;
ssl_ciphers ECDHE-RSA-AES256-GCM-SHA384:ECDHE-RSA-AES128-GCM-SHA256:
            ECDHE-RSA-CHACHA20-POLY1305:!CBC:!aNULL:!MD5:!DSS;
ssl_prefer_server_ciphers on;
\end{verbatim}
\textbf{Purpose:} Eliminates known TLS vulnerabilities (POODLE, BEAST, LUCKY13). The \texttt{!CBC} exclusion specifically mitigates LUCKY13 timing attacks. Only AEAD ciphers (GCM, CHACHA20-POLY1305) are permitted.\\
\textbf{Effectiveness:} \textcolor{shieldgreen}{\textbf{STRONG}} — Passed all 14 testssl.sh CVE checks.
\end{tcolorbox}

\subsection{HSTS (HTTP Strict Transport Security)}

\begin{tcolorbox}[colback=shieldgreen!5, colframe=shieldgreen]
\textbf{Mechanism:} \texttt{Strict-Transport-Security: max-age=31536000; includeSubDomains}\\
\textbf{Location:} \texttt{nginx/nginx.conf} (all location blocks)\\
\textbf{Purpose:} Instructs browsers to only connect via HTTPS for 1 year. Prevents SSL stripping attacks (sslstrip) by ensuring the browser never downgrades to HTTP after the first visit.\\
\textbf{Effectiveness:} \textcolor{shieldgreen}{\textbf{STRONG}} — 1 year max-age with subdomain coverage.
\end{tcolorbox}

\subsection{LAN-Only Access Restriction}

\begin{tcolorbox}[colback=shieldgreen!5, colframe=shieldgreen]
\textbf{Mechanism:} Nginx IP allowlist restricts access to internal networks only.\\
\textbf{Location:} \texttt{nginx/nginx.conf}\\
\textbf{Configuration:}
\begin{verbatim}
allow 192.168.0.0/16;
allow 10.0.0.0/8;
allow 127.0.0.1;
deny all;
\end{verbatim}
\textbf{Purpose:} The application is not internet-facing. Only IIIT-Delhi campus network clients can connect. This eliminates the entire class of internet-based attacks.\\
\textbf{Effectiveness:} \textcolor{shieldgreen}{\textbf{STRONG}} — External traffic is dropped at the Nginx layer.
\end{tcolorbox}

\subsection{BREACH Attack Prevention}

\begin{tcolorbox}[colback=shieldgreen!5, colframe=shieldgreen]
\textbf{Mechanism:} HTTP compression disabled (\texttt{gzip off}).\\
\textbf{Location:} \texttt{nginx/nginx.conf}\\
\textbf{Purpose:} The BREACH attack (CVE-2013-3587) exploits HTTP compression to extract secrets from encrypted HTTPS responses. By disabling gzip entirely, this attack vector is eliminated.\\
\textbf{Effectiveness:} \textcolor{shieldgreen}{\textbf{STRONG}} — Complete mitigation of CVE-2013-3587.
\end{tcolorbox}

% ================================================================
\section{Rate Limiting \& Connection Controls}
% ================================================================

\subsection{Dual-Layer Rate Limiting}

The platform implements rate limiting at both the reverse proxy and application layers:

\begin{center}
\small
\begin{tabularx}{\textwidth}{lXll}
\toprule
\textbf{Layer} & \textbf{Scope} & \textbf{Limit} & \textbf{Location} \\
\midrule
Nginx & Global API requests & 30 req/min per IP & \texttt{nginx.conf} \\
Nginx & Auth endpoints (login, register, TOTP, etc.) & 5 req/min per IP & \texttt{nginx.conf} \\
FastAPI & Global requests & 60 req/min per IP & \texttt{security.py} \\
FastAPI & Auth endpoints & 5 req/min per IP & \texttt{security.py} \\
\bottomrule
\end{tabularx}
\end{center}

\textbf{Why dual-layer?} Nginx rate limiting operates at the kernel level (efficient, low overhead) and blocks high-volume attacks before they reach the Python process. The application-level rate limiter provides finer-grained control over specific endpoint patterns and serves as a fallback if the Nginx layer is bypassed.

\subsection{Connection Throttling}

\begin{tcolorbox}[colback=shieldgreen!5, colframe=shieldgreen]
\textbf{Mechanism:} Maximum 15 concurrent connections per client IP.\\
\textbf{Location:} \texttt{nginx/nginx.conf} — \texttt{limit\_conn conn\_limit 15}\\
\textbf{Purpose:} Prevents Slowloris-type attacks where an attacker opens many connections and sends data slowly to exhaust server resources. Also limits the impact of connection flooding.\\
\textbf{Effectiveness:} \textcolor{shieldgreen}{\textbf{STRONG}} — Kernel-level enforcement.
\end{tcolorbox}

\subsection{Request Body Size Limits}

\begin{tcolorbox}[colback=shieldgreen!5, colframe=shieldgreen]
\textbf{Mechanism:} Request body size caps prevent large payload attacks.\\
\textbf{Location:} \texttt{nginx/nginx.conf}\\
\textbf{Limits:} 1 MB default (auth endpoints), 5 MB for API uploads.\\
\textbf{Purpose:} Prevents denial-of-service via large request bodies that could exhaust memory or disk space.\\
\textbf{Effectiveness:} \textcolor{shieldgreen}{\textbf{STRONG}} — Enforced at reverse proxy before reaching backend.
\end{tcolorbox}

% ================================================================
\section{Authentication \& Authorization Mitigation}
% ================================================================

\subsection{Argon2id Password Hashing}

\begin{tcolorbox}[colback=shieldgreen!5, colframe=shieldgreen]
\textbf{Mechanism:} Passwords are hashed using Argon2id (memory-hard, side-channel resistant).\\
\textbf{Location:} \texttt{backend/app/services/auth\_service.py}\\
\textbf{Purpose:} Argon2id is the winner of the Password Hashing Competition (PHC) and provides resistance against GPU-based brute-force attacks due to its memory-hard design. It also mitigates timing side-channel attacks.\\
\textbf{Effectiveness:} \textcolor{shieldgreen}{\textbf{STRONG}} — Industry-leading password hashing algorithm.
\end{tcolorbox}

\subsection{Password Policy}

\begin{tcolorbox}[colback=shieldgreen!5, colframe=shieldgreen]
\textbf{Mechanism:} Strict password requirements enforced via Pydantic validation.\\
\textbf{Location:} \texttt{backend/app/schemas/auth.py}\\
\textbf{Requirements:}
\begin{itemize}[nosep]
  \item Minimum 8 characters, maximum 72 characters
  \item At least one uppercase letter
  \item At least one lowercase letter
  \item At least one digit
  \item At least one special character
\end{itemize}
\textbf{Purpose:} Prevents weak passwords that are susceptible to dictionary attacks or brute-force.\\
\textbf{Effectiveness:} \textcolor{shieldgreen}{\textbf{STRONG}} — Enforced at schema level before any processing.
\end{tcolorbox}

\subsection{TOTP Two-Factor Authentication}

\begin{tcolorbox}[colback=shieldgreen!5, colframe=shieldgreen]
\textbf{Mechanism:} RFC 6238 TOTP mandatory for all users, with encrypted secret storage.\\
\textbf{Location:} \texttt{backend/app/services/totp\_service.py}\\
\textbf{Configuration:}
\begin{itemize}[nosep]
  \item Time step: 30 seconds
  \item Valid window: $\pm$1 step (90 seconds total)
  \item TOTP secrets encrypted with Fernet (AES-128-CBC + HMAC-SHA256) before storage
  \item QR code provided during setup; raw secret never returned in API responses
\end{itemize}
\textbf{Purpose:} Even if an attacker obtains a user's password, they cannot authenticate without the time-based code from the user's authenticator app.\\
\textbf{Effectiveness:} \textcolor{shieldgreen}{\textbf{STRONG}} — Mandatory 2FA with encrypted secret storage.
\end{tcolorbox}

\subsection{Progressive IP-Based Lockout}

\begin{tcolorbox}[colback=shieldgreen!5, colframe=shieldgreen]
\textbf{Mechanism:} Exponentially increasing lockout after failed TOTP attempts.\\
\textbf{Location:} \texttt{backend/app/services/totp\_service.py}\\
\textbf{Configuration:}
\begin{itemize}[nosep]
  \item 3 failed attempts within 5 minutes triggers lockout
  \item First lockout: 60 minutes
  \item Each subsequent lockout: $\times 5$ escalation (60 → 300 → 1500 → 7500 min)
  \item Keyed by client IP address (not email — prevents email-based DoS)
  \item Lockout state persists across server restarts (DB-backed)
\end{itemize}
\textbf{Purpose:} Makes brute-force attacks on TOTP codes computationally infeasible. After the first lockout, an attacker faces hours to days of waiting between attempts.\\
\textbf{Effectiveness:} \textcolor{shieldgreen}{\textbf{STRONG}} — Exponential backoff makes brute-force impractical.
\end{tcolorbox}

\subsection{JWT Token Security}

\begin{tcolorbox}[colback=shieldgreen!5, colframe=shieldgreen]
\textbf{Mechanism:} Short-lived access tokens with IP binding and type enforcement.\\
\textbf{Location:} \texttt{backend/app/services/auth\_service.py}\\
\textbf{Features:}
\begin{itemize}[nosep]
  \item Access token: 15-minute expiry
  \item Refresh token: 7-day expiry
  \item IP binding: token includes the client IP at issuance; requests from different IPs are rejected
  \item Type claim enforcement: refresh tokens cannot be used as access tokens
  \item Signed with HS256 using a cryptographically random 256-bit secret
\end{itemize}
\textbf{Purpose:} Short token lifetimes limit the window of exploitation for stolen tokens. IP binding prevents token reuse from different machines. Type enforcement prevents privilege escalation via token type confusion.\\
\textbf{Effectiveness:} \textcolor{shieldgreen}{\textbf{STRONG}} — Multi-layered token security.
\end{tcolorbox}

\subsection{Token Revocation}

\begin{tcolorbox}[colback=shieldgreen!5, colframe=shieldgreen]
\textbf{Mechanism:} SHA-256 hashed blocklist with DB persistence.\\
\textbf{Location:} \texttt{backend/app/services/auth\_service.py}\\
\textbf{Features:}
\begin{itemize}[nosep]
  \item Tokens are SHA-256 hashed before storage (raw JWTs never stored)
  \item In-memory cache for fast lookup; DB is source of truth
  \item Automatic cleanup of expired entries
  \item Revocation on logout, token refresh, and password change
  \item Persists across server restarts
\end{itemize}
\textbf{Purpose:} Allows immediate invalidation of compromised tokens, unlike pure JWT systems where tokens remain valid until expiry.\\
\textbf{Effectiveness:} \textcolor{shieldgreen}{\textbf{STRONG}} — Persistent revocation with O(1) lookup.
\end{tcolorbox}

% ================================================================
\section{Injection \& Input Validation Mitigation}
% ================================================================

\subsection{SQL Injection Prevention}

\begin{tcolorbox}[colback=shieldgreen!5, colframe=shieldgreen]
\textbf{Mechanism:} SQLAlchemy ORM with parameterized queries. LIKE/ILIKE wildcards escaped.\\
\textbf{Location:} All route files use SQLAlchemy ORM; \texttt{backend/app/utils/sanitize.py} provides \texttt{escape\_like()}.\\
\textbf{Purpose:} The ORM generates parameterized SQL statements, preventing SQL injection. The \texttt{escape\_like()} utility escapes \texttt{\%} and \texttt{\_} wildcards in user input used in LIKE/ILIKE clauses.\\
\textbf{Effectiveness:} \textcolor{shieldgreen}{\textbf{STRONG}} — No raw SQL anywhere in the codebase. Confirmed by sqlmap scan (0 injection points found).
\end{tcolorbox}

\subsection{XSS Prevention}

\begin{tcolorbox}[colback=shieldgreen!5, colframe=shieldgreen]
\textbf{Mechanism:} Multi-layered XSS defense.\\
\textbf{Controls:}
\begin{itemize}[nosep]
  \item \textbf{CSP:} \texttt{script-src 'self'} — no inline scripts, no external scripts
  \item \textbf{React:} Automatic DOM escaping of all rendered content
  \item \textbf{X-XSS-Protection:} \texttt{1; mode=block} — legacy browser protection
  \item \textbf{X-Content-Type-Options:} \texttt{nosniff} — prevents MIME-type sniffing
  \item \textbf{Pydantic validation:} Input length limits and type constraints
\end{itemize}
\textbf{Effectiveness:} \textcolor{shieldgreen}{\textbf{STRONG}} — Multiple overlapping controls.
\end{tcolorbox}

\subsection{CSRF Protection}

\begin{tcolorbox}[colback=shieldgreen!5, colframe=shieldgreen]
\textbf{Mechanism:} Cookie-based authentication with strict same-site policy.\\
\textbf{Location:} \texttt{backend/app/routes/auth.py}\\
\textbf{Controls:}
\begin{itemize}[nosep]
  \item \textbf{SameSite=strict:} Cookies are never sent on cross-origin requests
  \item \textbf{CORS:} Only same-origin requests are allowed
  \item \textbf{HttpOnly:} Cookies cannot be read by JavaScript (prevents exfiltration)
  \item \textbf{Secure:} Cookies only transmitted over HTTPS
\end{itemize}
\textbf{Purpose:} \texttt{SameSite=strict} is the strongest CSRF protection available — the browser will not include cookies in any request originating from a different site.\\
\textbf{Effectiveness:} \textcolor{shieldgreen}{\textbf{STRONG}} — Browser-enforced CSRF prevention.
\end{tcolorbox}

\subsection{CSWSH (Cross-Site WebSocket Hijacking) Protection}

\begin{tcolorbox}[colback=shieldgreen!5, colframe=shieldgreen]
\textbf{Mechanism:} WebSocket Origin header validation.\\
\textbf{Location:} \texttt{backend/app/routes/websocket.py}\\
\textbf{Implementation:} Before accepting a WebSocket connection, the server validates the \texttt{Origin} header against the allowed origin (\texttt{https://192.168.3.44}). Connections from unexpected origins are closed with code 4003.\\
\textbf{Purpose:} Prevents a malicious website from establishing a WebSocket connection to the backend using the user's authenticated cookies.\\
\textbf{Effectiveness:} \textcolor{shieldgreen}{\textbf{STRONG}} — Server-side origin enforcement.
\end{tcolorbox}

\subsection{File Upload Validation}

\begin{tcolorbox}[colback=shieldgreen!5, colframe=shieldgreen]
\textbf{Mechanism:} Magic bytes validation for all file uploads.\\
\textbf{Location:} \texttt{backend/app/services/resume\_service.py}, \texttt{backend/app/routes/users.py}\\
\textbf{Controls:}
\begin{itemize}[nosep]
  \item \textbf{Resumes:} PDF (\texttt{\%PDF}), DOCX (\texttt{PK}) magic bytes verified
  \item \textbf{Avatars:} JPEG (\texttt{0xFFD8FF}), PNG (\texttt{0x89504E47}), GIF (\texttt{GIF87a/89a}), WebP (\texttt{RIFF}+\texttt{WEBP}) verified
  \item File extension allowlist validation
  \item Request body size limits (5 MB for uploads)
\end{itemize}
\textbf{Purpose:} Prevents upload of executable files, HTML with XSS payload, or other malicious content disguised as allowed file types.\\
\textbf{Effectiveness:} \textcolor{shieldgreen}{\textbf{STRONG}} — Server-side magic bytes verification cannot be bypassed by MIME type spoofing.
\end{tcolorbox}

% ================================================================
\section{Security Headers}
% ================================================================

All security headers are set at the Nginx level with the \texttt{always} directive (applied to all response codes including errors). The backend maintains a defense-in-depth fallback.

\begin{center}
\small
\begin{tabularx}{\textwidth}{lXl}
\toprule
\textbf{Header} & \textbf{Value} & \textbf{Purpose} \\
\midrule
Strict-Transport-Security & max-age=31536000; includeSubDomains & Enforce HTTPS \\
X-Frame-Options & DENY & Prevent clickjacking \\
X-Content-Type-Options & nosniff & Prevent MIME sniffing \\
Referrer-Policy & strict-origin-when-cross-origin & Limit referrer leakage \\
Permissions-Policy & camera=(), microphone=(), geolocation=() & Disable browser features \\
X-XSS-Protection & 1; mode=block & Legacy XSS protection \\
Content-Security-Policy & default-src 'self'; script-src 'self'; style-src 'self' 'unsafe-inline'; img-src 'self' data: blob:; connect-src 'self'; font-src 'self'; object-src 'none'; frame-ancestors 'none'; base-uri 'self'; form-action 'self' & Comprehensive CSP \\
\bottomrule
\end{tabularx}
\end{center}

\textbf{CSP Notable Directives:}
\begin{itemize}[nosep]
  \item \texttt{object-src 'none'} — Blocks Flash and Java applets
  \item \texttt{frame-ancestors 'none'} — CSP-level clickjacking prevention (stronger than X-Frame-Options)
  \item \texttt{base-uri 'self'} — Prevents \texttt{<base>} tag injection for URL hijacking
  \item \texttt{form-action 'self'} — Prevents form submission to external domains
\end{itemize}

% ================================================================
\section{Path \& Information Disclosure Mitigation}
% ================================================================

\subsection{Sensitive Path Blocking}

Nginx returns 404 for all scanner-probed paths:

\begin{center}
\begin{tabular}{ll}
\toprule
\textbf{Pattern} & \textbf{Blocked Paths} \\
\midrule
Dotfiles & \texttt{/.env}, \texttt{/.git}, \texttt{/.svn}, \texttt{/.htaccess} \\
CMS probes & \texttt{/wp-admin}, \texttt{/wp-login} \\
Debug endpoints & \texttt{/phpinfo}, \texttt{/server-status}, \texttt{/server-info} \\
Frameworks & \texttt{/actuator}, \texttt{/console}, \texttt{/debug} \\
API docs & \texttt{/openapi.json}, \texttt{/api/docs}, \texttt{/api/redoc} \\
\bottomrule
\end{tabular}
\end{center}

\subsection{Information Concealment}

\begin{itemize}[nosep]
  \item \textbf{Server version:} \texttt{server\_tokens off} — hides Nginx version
  \item \textbf{Application version:} Removed from root endpoint response
  \item \textbf{Request ID:} Generated internally for audit trail but \textbf{not} exposed in response headers
  \item \textbf{API documentation:} Disabled when \texttt{DEBUG=False}
  \item \textbf{SQL queries:} Not echoed when \texttt{DEBUG=False}
\end{itemize}

% ================================================================
\section{Encryption at Rest}
% ================================================================

\begin{center}
\small
\begin{tabularx}{\textwidth}{lXl}
\toprule
\textbf{Data} & \textbf{Encryption} & \textbf{Location} \\
\midrule
Passwords & Argon2id hash (irreversible) & \texttt{auth\_service.py} \\
TOTP secrets & Fernet (AES-128-CBC + HMAC-SHA256) & \texttt{totp\_service.py} \\
Resumes & Fernet (AES-128-CBC + HMAC-SHA256) & \texttt{encryption\_service.py} \\
Messages & E2E RSA-4096 OAEP + AES-256-GCM & \texttt{pki\_service.py} \\
Server private keys & Passphrase-encrypted on disk & \texttt{pki\_service.py} \\
Audit trail & SHA-256 hash chaining & \texttt{audit\_service.py} \\
\bottomrule
\end{tabularx}
\end{center}

% ================================================================
\section{Mitigation Summary}
% ================================================================

\begin{center}
\begin{tabular}{clc}
\toprule
\textbf{\#} & \textbf{Mitigation Mechanism} & \textbf{Rating} \\
\midrule
1 & HTTPS enforcement (HTTP 444) & STRONG \\
2 & TLS 1.2/1.3 with AEAD ciphers only & STRONG \\
3 & HSTS (1-year max-age) & STRONG \\
4 & LAN-only access restriction & STRONG \\
5 & BREACH prevention (gzip off) & STRONG \\
6 & Nginx rate limiting (30/5 req/min) & STRONG \\
7 & Application rate limiting (60/5 req/min) & STRONG \\
8 & Connection throttling (15 per IP) & STRONG \\
9 & Request body size limits & STRONG \\
10 & Argon2id password hashing & STRONG \\
11 & Password complexity policy & STRONG \\
12 & TOTP 2FA (mandatory, encrypted secrets) & STRONG \\
13 & Progressive IP-based lockout & STRONG \\
14 & JWT security (IP binding, type enforcement) & STRONG \\
15 & Token revocation (DB-backed blocklist) & STRONG \\
16 & SQL injection prevention (ORM + escape\_like) & STRONG \\
17 & XSS prevention (CSP + React + headers) & STRONG \\
18 & CSRF protection (SameSite=strict) & STRONG \\
19 & CSWSH protection (Origin validation) & STRONG \\
20 & File upload validation (magic bytes) & STRONG \\
21 & Security headers (7 headers) & STRONG \\
22 & Path \& info disclosure blocking & STRONG \\
23 & Encryption at rest (6 data categories) & STRONG \\
\midrule
\multicolumn{2}{l}{\textbf{Total Mitigation Mechanisms}} & \textbf{23} \\
\bottomrule
\end{tabular}
\end{center}

\vspace{1em}

All 23 mitigation mechanisms are rated \textbf{STRONG}. The defense-in-depth strategy ensures that an attacker must bypass multiple independent controls to compromise any system component.

\end{document}
